% CV — Leshen Zhang (2026-09-23)
% Template: https://github.com/Yi-Lynn/cv-latex-template (FAANGPath Simple Template, resume.cls)
% Lato version: same content as cv_faangpath.tex, set in Lato.
\documentclass{resume} % Use the custom resume.cls style
\usepackage[left=0.4in,top=0.4in,right=0.4in,bottom=0.4in]{geometry} % Document margins
\usepackage[T1]{fontenc}
\usepackage[default]{lato} % Lato: open-source humanist sans, chosen for clarity + appearance
\usepackage{textcomp}
\usepackage{enumitem}
\usepackage{needspace}
\makeatletter\def\sectionskip{\smallskip}\makeatother
\setlist[itemize]{label=\textbullet,leftmargin=1.2em,itemsep=0.5pt,topsep=1pt,parsep=0pt,partopsep=0pt}
\newcommand{\me}{\textbf{\underline{Leshen Zhang}}}
\newcommand{\eq}{\textsuperscript{\textdaggerdbl}}
\name{Leshen Zhang}
\address{{\small \faEnvelope[regular]~\href{mailto:zhangls23@m.fudan.edu.cn}{zhangls23@m.fudan.edu.cn}}}
\address{{\small \faOrcid~\href{https://orcid.org/0009-0008-9143-6737}{0009-0008-9143-6737} \quad \faGlobe~\href{https://leshenzhang.github.io/}{leshenzhang.github.io}\quad \faGraduationCap~\href{https://scholar.google.com/citations?user=KiEU9eoAAAAJ}{Google Scholar}}}
\usepackage{fontawesome5}

\def\nameskip{\vspace{0pt}}\def\addressskip{\vspace{0pt}} % tighter header so two contact lines keep the CV at 3 pages
\begin{document}
\raggedright\hyphenpenalty=10000\exhyphenpenalty=10000 % ragged-right, no hyphenated words

\begin{rSection}{Education}
{\bf Fudan University} \hfill 09/2023 -- 07/2027 (expected) \\
Bachelor's Degree in Chemistry \textbar{} GPA: 3.88/4.00 \hfill \textit{Shanghai, China} \\
Advisor: Prof. Gengfeng Zheng
\end{rSection}

\begin{rSection}{Publications \hfill {\normalfont\footnotesize \textsuperscript{\textdaggerdbl}equal contribution, \textsuperscript{*}corresponding author}}
{\small\begin{itemize}[leftmargin=2.6em,labelsep=0.4em,itemsep=1pt]
  \item[\textbf{[J.1]}] \hypertarget{J.1}{}\me, Guangsheng~Liu, Xiao~Ma, Abdulrahman~Allangawi, Wan-Lu~Li*. \href{https://doi.org/10.1039/D6SC02460B}{\textbf{Defect-Mediated Regulation of Interfacial Hydrophobic Transport via Cavitation Thermodynamics}}. \\ \textit{Chem. Sci.} \textbf{2026}, 17, 15735--15744. \href{https://leshenzhang.github.io/papers/J1_Zhang2026_ChemSci.pdf}{[PDF]} \href{https://doi.org/10.1039/D6SC02460B}{[DOI]} (Sole First Author)
  \item[\textbf{[J.2]}] \hypertarget{J.2}{}Guangsheng~Liu\eq, Xiao~Ma\eq, \me\eq, Jenedith~Pascasio, Jonathan~Yang, Yuxiang~Chen, Wan-Lu~Li*. \href{https://doi.org/10.1039/D6DD00273K}{\textbf{CatGo: Bridging CLI Coding Agents with Interactive Structure and Workflow Management for Computational Chemistry}}. \\ \textit{Digit. Discov.} \textbf{2026}, 5, 3825--3849. \href{https://leshenzhang.github.io/papers/J2_Liu2026_DigitDiscov.pdf}{[PDF]} \href{https://doi.org/10.1039/D6DD00273K}{[DOI]} (Co-First Author, Rank Third)
  \item[\textbf{[J.3]}] \hypertarget{J.3}{}Shuya~Hao\eq, Maoyin~Wang\eq, \me, Ximeng~Lv, Chen~Peng, Yuhang~Huang, Pei~Yuan*, Qing~Han*, Gengfeng~Zheng*. \href{https://doi.org/10.1002/anie.202510241}{\textbf{Switching Photocatalytic Methane Oxidation Toward Ethanol by Tuning Spin States}}. \\ \textit{Angew. Chem. Int. Ed.} \textbf{2025}, 64, e202510241. \href{https://leshenzhang.github.io/papers/J3_Hao2025_AngewChem.pdf}{[PDF]} \href{https://doi.org/10.1002/anie.202510241}{[DOI]} (Third Author)
\end{itemize}}
\end{rSection}

\begin{rSection}{Preprints \& Papers Under Review \hfill {\normalfont\footnotesize \textsuperscript{\textdaggerdbl}equal contribution, \textsuperscript{*}corresponding author}}

{\small\begin{itemize}[leftmargin=2.6em,labelsep=0.4em,itemsep=1pt]
  \item[\textbf{[S.1]}] \hypertarget{S.1}{}\me\eq, Xujian~Wang\eq, Taoyu~Niu, Xuchao~Zhou, Xiao~Ma, Guangsheng~Liu, Olexandr~Isayev*, Junmei~Wang*, Wan-Lu~Li*. \textbf{Conformation-Responsive ML/MM Charges in PBSA for Protein--Ligand Affinity Ranking}. \\ To be submitted. (Co-First Author, Rank First)
  \item[\textbf{[S.2]}] \hypertarget{S.2}{}Guangsheng~Liu\eq, \me\eq, Xiao~Ma, Abdulrahman~Allangawi, Huabin~Zhang*, Wan-Lu~Li*. \textbf{Local-Environment Populations Govern Size-Dependent OER Activity in Supported Subnanometer IrO\textsubscript{x} Clusters}. \\ \textit{Adv. Funct. Mater.}, under review. (Co-First Author, Rank Second)
  \item[\textbf{[S.3]}] \hypertarget{S.3}{}Xiao~Ma, \me, Guangsheng~Liu, Abdulrahman~Allangawi, Huabin~Zhang*, Wan-Lu~Li*. \href{https://doi.org/10.26434/chemrxiv.15008697/v1}{\textbf{Surface Geometry Governs Cation Effects via Nonlocal Solvent Gating}}. \\ \textit{Angew. Chem. Int. Ed.}, under review. \href{https://chemrxiv.org/doi/pdf/10.26434/chemrxiv.15008697/v1}{[PDF]} \href{https://doi.org/10.26434/chemrxiv.15008697/v1}{[DOI]} (Second Author)
  \item[\textbf{[S.4]}] \hypertarget{S.4}{}Guangsheng~Liu\eq, Xiao~Ma\eq, \me\eq, Xujian~Wang, Abdulrahman~Allangawi, Shaowei~Li, Wan-Lu~Li*. \href{https://doi.org/10.26434/chemrxiv.15003924/v1}{\textbf{Intercalation-Modulated Electronic Regulation of Nitrogen Reduction on Doped MoS\textsubscript{2}}}. \\ To be submitted. \href{https://leshenzhang.github.io/papers/S4_Liu2026_ChemRxiv.pdf}{[PDF]} \href{https://doi.org/10.26434/chemrxiv.15003924/v1}{[DOI]} (Co-First Author, Rank Third)
  \item[\textbf{[S.5]}] \hypertarget{S.5}{}Guangsheng~Liu\eq, \me\eq, Xiao~Ma, I-Shan~Tsai, Jinyi~Zhang, Wan-Lu~Li*. \href{https://doi.org/10.26434/chemrxiv.15006055/v1}{\textbf{ChemMate: An Autonomous Research Mate for the Full Computational Chemistry Research Lifecycle, from Ideation to Manuscript}}. \\ To be submitted. \href{https://leshenzhang.github.io/papers/S5_Liu2026_ChemRxiv.pdf}{[PDF]} \href{https://doi.org/10.26434/chemrxiv.15006055/v1}{[DOI]} (Co-First Author, Rank Second)
  \item[\textbf{[S.6]}] \hypertarget{S.6}{}Tingsong~Li\eq, Jifeng~Wu\eq, Abdulrahman~Allangawi\eq, \me, Wan-Lu~Li*, Xiangyun~Xiao, Miao~Hu, Mohd~Adnan~Khan, Rashed~Aleisa, Shouwei~Zuo, Heng~Guo, Kuo-Wei~Huang*, Ying~Zhou*, Huabin~Zhang*, Magnus~Rueping*. \textbf{Interfacial Proton Relay Enables CO\textsubscript{2}-to-Methanol Electroreduction on Layered Cobalt Phthalocyanine}. \\ \textit{J. Am. Chem. Soc.}, under review. (Fourth Author)
  \item[\textbf{[S.7]}] \hypertarget{S.7}{}Xuchao~Zhou, Xujian~Wang, Jiahao~Xie, \me, Junmei~Wang*. \href{https://doi.org/10.26434/chemrxiv.15008812/v1}{\textbf{Leveraging Machine Learning Interatomic Potentials into Fitting Fast and Accurate AMBER-Compatible Precise Force Field}}. \\ \textit{Chem. Sci.}, under review. \href{https://chemrxiv.org/doi/pdf/10.26434/chemrxiv.15008812/v1}{[PDF]} \href{https://doi.org/10.26434/chemrxiv.15008812/v1}{[DOI]} (Fourth Author)
  \item[\textbf{[S.8]}] \hypertarget{S.8}{}Yunzhi~Wang\eq, Chengyang~Feng\eq, Lijie~Wang\eq, Yang~Cheng, \me, Yinglu~Jia, Lingyun~Zhao, Abdulrahman~Allangawi, Weiwei~Li*, Wan-Lu~Li*, Martin~Heeney, Magnus~Rueping, Omar~F.~Mohammed*, Huabin~Zhang*. \textbf{Overcoming Excitonic Bottlenecks via Molecular Intercalation into Donor Domains for NIR-Driven Hydrogen Evolution}. \\ \textit{Nat. Commun.}, under review. (Fifth Author)
\end{itemize}}
\end{rSection}

\Needspace{7\baselineskip}
\begin{rSection}{Selected Research Experience}
\interlinepenalty=10000 % never break a project entry across pages
\Needspace{6\baselineskip}\textbf{I. First-Principles Calculations for Materials and Catalysis}\par
\Needspace{5\baselineskip}{\bfseries\itshape (a) Surface Reconstruction and Interfacial Structure}\par
\begin{itemize}[itemsep=6pt]
  \item \textbf{Defect-Mediated Regulation of Interfacial Hydrophobic Transport via Cavitation Thermodynamics} \\* {\small\hyperlink{J.1}{\textbf{[J.1]}}}\hfill{\small\textit{PI: \textbf{Prof. Wan-Lu Li}, UC San Diego, USA}} \\* Revealed that surface defects regulate interfacial hydrophobic transport by tuning the local cavitation free energy, establishing cavitation thermodynamics as a descriptor linking interfacial water structure to reactivity
  \item \textbf{Local-Environment Populations Govern Size-Dependent OER Activity in Supported Subnanometer IrO\textsubscript{x} Clusters} \\* {\small\hyperlink{S.2}{\textbf{[S.2]}}}\hfill{\small\textit{PI: \textbf{Prof. Wan-Lu Li}, UC San Diego, USA}} \\* Formulated a constant-potential ensemble-sampling methodology that maps site-specific overpotentials and activation barriers across cluster sizes and configurations, bridging local site motifs with microkinetics to identify the ideal size window for Ir clusters
  \item \textbf{Surface Geometry Governs Cation Effects via Nonlocal Solvent Gating} \\* {\small\hyperlink{S.3}{\textbf{[S.3]}}}\hfill{\small\textit{PI: \textbf{Prof. Wan-Lu Li}, UC San Diego, USA}} \\* Showed that Au surface geometry governs alkali-cation effects in CO\textsubscript{2} reduction: interfacial water reverses the cation effects and propagates them nonlocally to pristine terraces, explaining the volcano-shaped cation dependence on polycrystalline Au
  \item \textbf{Site-Resolved C--C Coupling on Restructured Cu Electrodes (ongoing)} \\* \hspace*{\fill}{\small\textit{PI: \textbf{Prof. Anastassia N. Alexandrova}, UC Los Angeles, USA}} \\* Investigating how kinetically trapped, adsorbate-roughened Cu electrodes catalyze C--C coupling during CO\textsubscript{2} reduction, and using site-resolved KMC simulation to identify the classes of restructured Cu sites that dominate the coupling rate
\end{itemize}
\Needspace{6\baselineskip}{\bfseries\itshape (b) Catalytic Mechanisms and Multiscale Modeling}\par
\begin{itemize}[itemsep=6pt]
  \item \textbf{Switching Photocatalytic Methane Oxidation Toward Ethanol by Tuning Spin States} \\* {\small\hyperlink{J.3}{\textbf{[J.3]}}}\hfill{\small\textit{PI: \textbf{Prof. Gengfeng Zheng}, Fudan University, China}} \\* \textit{Experiment:} developed a Zn--O--Fe catalyst that selectively photo-oxidizes CH\textsubscript{4} to ethanol \\ \textit{Computation:} revealed that the medium Fe spin state preferentially stabilizes \textbullet{}CH\textsubscript{2}OH over \textbullet{}OCH\textsubscript{3}, explaining why the catalyst with medium-spin polarization favors ethanol
  \item \textbf{Intercalation-Modulated Electronic Regulation of Nitrogen Reduction on Doped MoS\textsubscript{2}} \\* {\small\hyperlink{S.4}{\textbf{[S.4]}}}\hfill{\small\textit{PI: \textbf{Prof. Wan-Lu Li}, UC San Diego, USA}} \\* Developed a screening strategy that excludes unstable catalysts before DFT and distilled an interpretable SISSO descriptor governing NRR activity and selectivity against HER
  \item \textbf{Multiscale Modeling of PEM Water-Electrolysis OER: From DFT to Device} \\* \hspace*{\fill}{\small\textit{PI: \textbf{Prof. Wan-Lu Li}, UC San Diego, USA}} \\* Bridged DFT, microkinetic modeling, and mass-transfer simulation in a multiscale model, showing that catalyst morphology governs local mass transport and thereby the apparent activity
\end{itemize}
\Needspace{7\baselineskip}\textbf{II. AI-Assisted Scientific Software}\par
\begin{itemize}[itemsep=6pt]
  \item \textbf{CatGo: Bridging CLI Coding Agents with Interactive Structure and Workflow Management for Computational Chemistry} \\* {\small\hyperlink{J.2}{\textbf{[J.2]}}}\hfill{\small\textit{PI: \textbf{Prof. Wan-Lu Li}, UC San Diego, USA}} \\* Established an open-source platform unifying the structure \textrightarrow{} DFT \textrightarrow{} HPC \textrightarrow{} analysis chain, and bridged AI coding agents to its workbench, enabling high-throughput calculations
  \item \textbf{ChemMate: A Conversational Research Collaborator for Computational Chemistry} \\* {\small\hyperlink{S.5}{\textbf{[S.5]}}}\hfill{\small\textit{PI: \textbf{Prof. Wan-Lu Li}, UC San Diego, USA}} \\* Established an autonomous research collaborator that runs entire projects from ideation to manuscript under human gates, demonstrated in five campaigns with reproducible results
\end{itemize}
\Needspace{7\baselineskip}\textbf{III. Molecular Modeling: Electrostatics in Enzyme Design}\par
\begin{itemize}[itemsep=6pt]
  \item \textbf{Conformation-Responsive ML/MM Charges in PBSA for Protein--Ligand Affinity Ranking} \\* {\small\hyperlink{S.1}{\textbf{[S.1]}}}\hfill{\small\textit{PI: \textbf{Prof. Wan-Lu Li}, UC San Diego, and \textbf{Prof. Junmei Wang}, University of Pittsburgh, USA}} \\* Developed the PCME method for ML/MM embedding, which feeds conformation-resolved ML charges of the ligand and binding site into the Poisson--Boltzmann solver, accurately predicting binding affinities with high computational efficiency
  \item \textbf{Leveraging Machine Learning Interatomic Potentials into Fitting Fast and Accurate AMBER-Compatible Precise Force Field} \\* {\small\hyperlink{S.7}{\textbf{[S.7]}}}\hfill{\small\textit{PI: \textbf{Prof. Junmei Wang}, University of Pittsburgh, USA}} \\* Established an MLIP-driven workflow that builds AMBER-compatible force fields, matching QM-refined accuracy in hydration free energies at far lower cost and extending to noncanonical residues and metal sites in enzymes

\end{itemize}
\end{rSection}

\begin{rSection}{Skills}
\textbf{Experimental Skills}
\begin{itemize}
  \item Inorganic/organic synthesis; characterization by GC, HPLC, UV-Vis, FT-IR, and NMR
  \item Electro- and photocatalytic testing with H-cells, flow cells, and membrane electrode assemblies
\end{itemize}
\textbf{Theoretical Chemistry Skills}
\begin{itemize}
  \item \textbf{First principles:} DFT, GC-DFT, AIMD, and enhanced sampling (VASP, CP2K, Gaussian)
  \item \textbf{Molecular simulation:} ML/MM, ML-MD (Amber)
\end{itemize}
\textbf{Machine Learning Skills}
\begin{itemize}
  \item \textbf{Models:} graph, convolutional, and recurrent neural networks; ML interatomic potentials
  \item \textbf{Applications:} MLIP training, interpretable descriptor discovery, and molecular classification
\end{itemize}
\textbf{AI-Assisted Software Development}
\begin{itemize}
  \item \textbf{Research software:} co-developed CatGo (GitHub: \href{https://github.com/Hello-QM/catgo-LRG}{Hello-QM/catgo-LRG}), ChemMate, and agent-workspace plugins, architected and built with AI coding agents; independently developed SPARK (GitHub: \href{https://github.com/WanluLigroupUCSD/SPARK}{WanluLigroupUCSD/SPARK}), an on-the-fly kinetic Monte Carlo (KMC) package
  \item \textbf{Agent design:} human-in-the-loop agent tools (preview \textrightarrow{} approval \textrightarrow{} apply)
\end{itemize}
\end{rSection}

\begin{rSection}{Honors and Awards}
\begin{itemize}
  \item NSFC Grant --- Pilot Program for Young Students Basic Research (Undergraduate), National Natural Science Foundation of China \hfill 09/2026 -- 07/2027
  \item Fudan University `Syensqo' Scholarship \hfill 11/2025
  \item Fudan University `Ruiqing' Scholarship \hfill 11/2024
  \item Fudan University \& Chemours `Future of Chemistry' Scholarship \hfill 11/2024
  \item Second Prize, 36th China Chemistry Olympiad (National Final) \hfill 11/2022
\end{itemize}
\end{rSection}

\begin{rSection}{Scientific Presentations}
\Needspace{4\baselineskip}\textbf{Oral Presentations}
\begin{itemize}
  \item ``Defect-Mediated Regulation of Hydrophobic Transport via Cavitation Thermodynamics'' \hfill \mbox{12/2026 (expected)} \\
  \textit{Chinese Chemical Society, The 16th National Conference of Theoretical and Computational Chemistry}, Shenzhen, China
  \item ``Spin Polarization in Photocatalysis'' \hfill 11/2025 \\
  \textit{Fudan University, Undergraduate Student Forum 2025 on Chemistry}, Shanghai, China
\end{itemize}
\Needspace{4\baselineskip}\textbf{Poster Presentations}
\begin{itemize}
  \item ``Switching Photocatalytic Methane Oxidation Toward Ethanol by Tuning Spin States'' \hfill 11/2025 \\
  \textit{University of Chinese Academy of Sciences, Graduate Student Forum 2025 on Energy Catalysis}, Dalian, China
\end{itemize}
\end{rSection}
\Needspace{3\baselineskip}
\begin{rSection}{Professional Service}
\textbf{Peer Reviewer}, \textit{J. Colloid Interface Sci.}, 5 completed reviews \hfill 2024 -- 2026
\end{rSection}

\end{document}
